\chapter{Herramientas, requerimientos de software y arquitectura}
\label{chp:4}


Para dar cumplimiento al requerimiento de uso de software libre, el cuadro~\ref{tab:software} presenta el conjunto de programas empleados durante el desarrollo del proyecto, junto con una breve descripción de su uso y la licencia correspondiente. Para la clasificación de las licencias se adopta la definición de \ac{FLOSS}, que incluye software libre y de código abierto, según lo establecido por la \ac{FSF}\cite{gnu_free_sw}.

El único software utilizado que no se enmarca dentro de la categoría \acs{FLOSS} es GitHub, el cual fue empleado como plataforma de alojamiento del código fuente y como herramienta de colaboración entre los integrantes del equipo. No obstante, el proyecto no presenta \textit{vendor lock-in}, ya que la dependencia de esta plataforma no es estructural. En consecuencia, es posible —y se recomienda— migrar el repositorio y los flujos de trabajo a una instancia propia de Gitea\cite{gitea_about}, lo que permitiría que el proyecto sea completamente desarrollado y gestionado con software libre y abierto.

\begin{longtable}{|p{0.15\textwidth}|p{0.4\textwidth}|p{0.3\textwidth}|}
  \caption{Software utilizado, descripción y licencias}
  \label{tab:software} \\
  \hline
  \textbf{Software} & \textbf{Descripción} & \textbf{Licencia} \\ 
  \hline
  \endfirsthead

  \hline
  \textbf{Software} & \textbf{Descripción} & \textbf{Licencia} \\ 
  \hline
  \endhead

  \hline
  \multicolumn{3}{r}{\emph{Continúa en la siguiente página}} \\
  \endfoot

  \hline
  \endlastfoot

  Python & Lenguaje de programación usado durante todo el proyecto 
                                           & PSFL (Open Source)\cite{python_license} \\ \hline

  Docker & Software de virtualización para ayudar al despliegue de todos los servicios 
                                           & Apache 2.0 (Open Source)\cite{moby_license} \\ \hline

  PostGIS & Modulo para soporte de objetos geográficos a PostgreSQL 
                                           & GPL-2 (Libre)\cite{postgis} \\ \hline

  NGINX & Usado como servidor web dentro del proyecto 
                                           & BSD-2 (Open Source)\cite{nginx_license} \\ \hline

  GitHub & Servicio de hosting y colaboración para alojar el código 
                                           & Propietario\cite{wiki_github} \\ \hline

  QGIS & Análisis de datos geoespaciales 
                                           & GPL-2 (Libre)\cite{qgis_license} \\ \hline

  Marimo & Notebooks para trabajar, limpiar y procesar datos 
                                           & Apache 2.0 (Open Source)\cite{marimo_license} \\ \hline

  Svelte & Framework de \acs{JS} usado para el front-end del proyecto 
                                           & MIT (Open Source)\cite{svelte_license} \\ \hline

  Leaflet & Librería de \acs{JS} para mostrar y montar capas de datos geográficos teniendo como capa base OSM 
                                           & BSD-2 (Open Source)\cite{leaflet_license} \\ \hline

  FastAPI & Framework de Python utilizado para desarrollar las \acs{API} del proyecto 
                                           & MIT (Open Source)\cite{fastapi_license} \\ \hline

  Drawio & Aplicación utilizada para dibujar la arquitectura & Apache 2.0 (Open Source)\cite{drawio_license} \\ \hline

  GNU Emacs & Editor de texto utilizado para editar archivos de datos, usado por sus macros & GPL-3 (Libre)\cite{wiki_gnu_emacs} \\ \hline

  \LaTeX & Sistema de software para la creación del presente documento & LPPL (Libre)\cite{latex_lppl} \\ \hline

  Bash & Lenguaje de scripting usado para automatizar scripts & GPL-3 (Libre)\cite{wiki_bash_shell} \\ \hline

  GNU Make & Utilidad usada para compilar algunos archivos del proyecto & GPL-3 (Libre)\cite{gnu_make_manual} \\ \hline

  LibreOffice Calc & Hojas de cálculo para realizar operaciones con datos & MPL-2 (Open Source)\cite{libreoffice_calc} \\ \hline

  Git & Software de control de versiones distribuido para el desarrollo de todo el proyecto & GPL-2 (Libre)\cite{git_about} \\ \hline 
  
\end{longtable}



